% 第四节：齐次坐标
\chapter{齐次坐标}

\section{圆锥曲线的一般方程}

以椭圆为例，当其中心在原点、长轴与 $x$ 轴重合时，标准方程为
\[
\frac{x^2}{a^2}+\frac{y^2}{b^2}=1 .
\]
但椭圆可以有任意的位置与姿态，此时它的方程的一般形式为
\begin{equation}\label{eq:gen}
ax^2+2bxy+cy^2+2dx+2ey+f=0 .
\end{equation}
在系数取适当值时，这个方程可以表示椭圆、抛物线或双曲线；在系数取某些特殊值时，它也可能表示两条直线、一条直线或一个点，或者无解。以下先讨论一般情形，特殊情形暂不涉及。

为了消去常数项与一次项，作代换
\[
x=\frac{X}{Z},\qquad y=\frac{Y}{Z},\qquad Z\neq 0,
\]
代入 \eqref{eq:gen} 并乘以 $Z^2$，得到
\begin{equation}\label{eq:hom}
aX^2+2bXY+cY^2+2dXZ+2eYZ+fZ^2=0 .
\end{equation}
这个方程关于 $X,Y,Z$ 的每一项都是二次的。进一步引入对称的 $(0,2)$ 型张量 $g$，其分量为
\[
g_{11}=a,\quad g_{12}=g_{21}=b,\quad g_{22}=c,\quad
g_{13}=g_{31}=d,\quad g_{23}=g_{32}=e,\quad g_{33}=f,
\]
并把未知量 $X,Y,Z$ 看作矢量 $A$ 的分量 $A^1=X$、$A^2=Y$、$A^3=Z$，则 \eqref{eq:hom} 可改写为
\begin{equation}\label{eq:gAA}
g_{ij}A^iA^j=0 .
\end{equation}
若把 $g$ 取作度规并用它定义点乘，则 \eqref{eq:gAA} 可进一步简写为
\[
A^2=g_{ij}A^iA^j=0 .
\]
也就是说，圆锥曲线可以用“矢量平方为零”这一与坐标系无关的方程来描述。正是这一简化的形式提示了把平面几何的问题放到三维矢量空间中处理的思路。

\begin{remark}
方程 $ax^2+2bxy+cy^2+2dx+2ey+f=0$ 在系数退化时可能表示两条直线、一条直线、一个点，或者没有实点。这些情形对应 $g$ 的退化或方程的实解集退化，本书后文在需要时分别处理，本节只讨论 $g$ 非退化的一般情形。
\end{remark}

\section{平面上的点与三维矢量的对应}

上节的改写表明，平面上的点可以借助三个分量 $X,Y,Z$ 来表示。下面把这个对应关系明确地建立起来。整个过程分三步。

第一步：在三维矢量空间中建立坐标系，把任意矢量写成
\[
A^i=(X,Y,Z).
\]

第二步：利用这个坐标系，把平面上的点 $(x,y)$ 与某个确定方向上的全体非零三维矢量对应起来。约定的对应关系是：点 $(x,y)$ 对应所有分量形如
\[
A^i=(\lambda x,\lambda y,\lambda),\qquad \lambda\in\R\setminus\{0\}
\]
的三维矢量；反之，矢量 $(X,Y,Z)$ 以及任何形如 $(\lambda X,\lambda Y,\lambda Z)$ 的矢量都对应平面上的点
\[
\left(\frac{X}{Z},\frac{Y}{Z}\right).
\]
这个对应是多对一的：平面上的每个点对应一条过原点的直线（去掉原点）上的全部非零矢量，而每个这样的矢量只对应一个点。

\begin{definition}[齐次坐标]
按上述约定，与平面上的点 $(x,y)$ 相对应的三维矢量的分量 $(\lambda x,\lambda y,\lambda)$（$\lambda\neq 0$）称为该点的一组齐次坐标；反过来，矢量 $(X,Y,Z)$（$Z\neq 0$）对应的点是 $(X/Z,Y/Z)$。
\end{definition}

当 $Z=0$ 时上述对应没有定义，因为分母为零，此时矢量不对应任何平面上的点。这一类矢量所对应的几何对象将在讨论无穷远时处理。

第三步：抹去前两步中建立坐标系所留下的一切痕迹。

在完成前两步之后，平面上的几何对象及其相互关系都可以用带有坐标系的三维矢量表示。具体地说，存在两个函数 $x(A),y(A)$，它们给出矢量 $A$ 所对应点的坐标；而平面几何的关系总可以写成关于 $x,y$ 的方程，因此也就可以改写成关于 $A$ 的方程。例如曲线 $y=\sin x$ 等价于关于 $A$ 的方程
\[
y(A)=\sin x(A).
\]
然而，对这样的曲线使用三维矢量的方法求解没有任何好处：求解这个方程终究要求出矢量的分量，再让分量满足 $x(A),y(A)$ 的复杂关系，这层包装属于画蛇添足。

三维矢量方法的优势只出现在下述情形：关于矢量的方程由简单自然的运算组成，其中尤其重要的是与坐标系无关的运算，例如加法、数乘，以及定义了度规之后的点乘。这就引出下面的定义。

\begin{definition}[射影几何学]
如果平面几何量之间的关系经过前两步的对应改写为关于三维矢量的等式之后，所得的等式与坐标系无关，则称这个几何关系属于射影几何学。
\end{definition}

对于射影几何学的问题，用三维矢量方法求解是最简便的做法。换言之，只有当第三步（抹去坐标系的一切痕迹）可以做到时，使用三维矢量方法才具有优势；若某个几何关系不得不借助坐标系来表达，则它不属于射影几何学。
